\subsection{Hipótesis externa}

El fuerte Moran's $I$ y la sensibilidad municipal sugieren que una estadística pública
localizada podría aportar un prior regional. SIAP 2025 es especialmente interesante porque es
contemporáneo y público. La hipótesis de 04E.1 no es que SIAP sea equivalente al target
parcelario, sino que
\[
y_i = g(z_{m(i),2025}) + r_i,
\]
donde $z_{m(i),2025}$ es una señal municipal y $r_i$ una corrección parcelaria, podría tener
menor varianza que modelar $y_i$ sin prior.

Para probarla de manera semánticamente correcta se vuelve a la tabla SIAP detallada, evitando
la agregación legada que mezclaba ciclos/modalidades.

\subsection{Scope exacto y fallbacks explícitos}

El scope principal se define como:
\[
\text{Cebada grano}
+\text{Primavera--Verano}
+\text{Temporal}
+\text{CVEGEO}
+\text{2025}.
\]
La elección de Primavera--Verano se alinea con el ciclo abril--octubre del reto. En este
documento se mantiene como decisión operacional del proyecto, no como una validación externa
independiente de nomenclatura SIAP.

La auditoría conserva scopes progresivamente más amplios:
\begin{description}
\item[exact] ciclo y modalidad exactos;
\item[cycle all-mode] ciclo exacto, modalidad libre;
\item[temporal all-cycle] modalidad Temporal, ciclo libre;
\item[all-grain] toda Cebada grano.
\end{description}
Los fallbacks se nombran en el artefacto; nunca se mezclan silenciosamente.

\begin{table}[H]
\centering
\caption{Auditoría de scopes SIAP en 04E.1.}
\label{tab:cp04e-scope}
\begin{tabular}{lrrr}
\toprule
Scope & filas históricas & municipios 2025 & filas 2025 \\
\midrule
exact & 924 & 83 & 83 \\
cycle all-mode & 925 & 83 & 93 \\
temporal all-cycle & 929 & 83 & 83 \\
all-grain & 930 & 83 & 94 \\
\bottomrule
\end{tabular}
\end{table}

La cobertura exacta alcanza 197/197 parcelas. El prior seleccionado tiene cobertura 197/197
antes de fallback por mediana, incluyendo 59/59 objetivos reales. Por tanto, cualquier fracaso
predictivo no puede atribuirse a missingness de SIAP.

\safefigure{../../reports/checkpoint_04e1/figures/siap_scope_coverage.png}
{Cobertura de los scopes SIAP auditados.}
{fig:cp04e-coverage}

\subsection{Construcción del rendimiento municipal}

En un municipio-año, la agregación evita promediar rendimientos reportados sin peso cuando
existen varias filas. Se conservan tanto el rendimiento reconstruido
\[
Y^{prod/harv}_{m,t}
=
\frac{\sum_r \mathrm{VolumenProduccion}_{r}}
{\sum_r \mathrm{Cosechada}_{r}},
\]
como el rendimiento reportado ponderado por superficie cuando corresponde. También se deriva
tasa de siniestro. El objetivo es mantener una definición reproducible aun cuando el source
tenga varias categorías administrativas.

\subsection{Usos evaluados del prior}

Se probaron cuatro formas principales.

\paragraph{Directo.}
\[
\hat y_i=z_{m(i),2025}.
\]

\paragraph{Calibración affine.}
Sobre etiquetas visibles:
\[
y_i=\beta_0+\beta_1z_i+\varepsilon_i,
\]
con Ridge sobre $(1,z_i)$.

\paragraph{Local residual.}
Se define
\[
r_i=y_i-z_i
\]
para etiquetas visibles y se ajusta Local Ridge para $r_q$:
\[
\hat y_q=z_q+\hat r_q.
\]

\paragraph{Graph residual.}
Los residuos visibles se propagan sobre el grafo:
\[
\hat y_q=z_q+\operatorname{Graph}(r)_q.
\]

Finalmente, se prueban blends fijos del 25, 50 y 75\% contra los baselines 04D.

\subsection{Resultados target-matched}

\begin{table}[H]
\centering
\caption{Resultados principales de SIAP 04E.1.}
\label{tab:cp04e-results}
\small
\begin{tabular}{lrr}
\toprule
Método & RMSE medio & pooled MAE \\
\midrule
Local04D & \textbf{0.4878} & 0.3425 \\
Local + SIAP-local 25\% & 0.4922 & 0.3528 \\
Graph04D & 0.4949 & \textbf{0.3403} \\
Graph + SIAP-graph 25\% & 0.5014 & 0.3500 \\
SIAP local residual & 0.5456 & 0.4151 \\
SIAP graph residual & 0.5558 & 0.3976 \\
SIAP affine & $\approx0.837$ & $\approx0.717$ \\
SIAP directo & 1.4368 & 1.3095 \\
\bottomrule
\end{tabular}
\end{table}

La cobertura perfecta no se traduce en resolución parcelaria. En las 138 etiquetas completas:
\[
\RMSE(z,y)=1.4926,
\qquad
\rho_{\mathrm{Spearman}}(z,y)=-0.2181.
\]
La feature SIAP 2025 legada produce RMSE directo 1.4894, prácticamente igual de pobre. Filtrar
correctamente ciclo/modalidad mejora semántica, no garantiza señal parcelaria.

\safefigure{../../reports/checkpoint_04e1/figures/siap_proxy_vs_yield.png}
{Proxy SIAP municipal frente al rendimiento observado de parcela.}
{fig:cp04e-siap-yield}

\subsection{LOSO y evidencia contra la incorporación}

El selector leave-one-target-matched-split-out dispone de Local04D, Graph04D y las variantes
SIAP. En 16/16 holdouts elige Local04D. Esta unanimidad es evidencia más fuerte que comparar
el mejor promedio de una tabla completa: cada holdout se evalúa después de seleccionar con los
otros splits.

\subsection{Resultado grouped que sí favorece SIAP}

El fracaso primario no es uniforme. Bajo municipality-grouped:
\[
\RMSE(\mathrm{Graph04D})=0.5279,
\]
mientras un blend 25\% Graph/SIAP llega a aproximadamente
\[
0.5096.
\]
Esto tiene interpretación coherente: cuando se elimina un municipio completo, un prior
municipal externo puede suplir parte de la información que ya no existe en los labels
cercanos.

Sin embargo, el test real no es un leave-one-municipality-out puro. Las máscaras
target-matched son la simulación diseñada para la geometría específica de los 59 objetivos, y
allí el mismo blend empeora. Por ello el resultado grouped se conserva como stress evidence,
no como autorización para introducir SIAP al predictor final.

\safefigure{../../reports/checkpoint_04e1/figures/external_method_ranking.png}
{Ranking de métodos con evidencia SIAP en el protocolo target-matched.}
{fig:cp04e-ranking}

\subsection{Lección de resolución espacial}

Este checkpoint ilustra un principio general de data fusion: la contemporaneidad no compensa
una discordancia de resolución. SIAP observa un outcome administrativo municipal; el target es
parcelario. Si dentro de un municipio existe heterogeneidad considerable, el prior
$z_m$ tiene un error irreducible para parcelas individuales.

La evidencia externa sigue siendo metodológicamente valiosa porque:
\begin{itemize}
\item verifica una hipótesis razonable;
\item documenta el significado exacto de SIAP;
\item cuantifica su utilidad bajo extrapolación municipal;
\item evita que el equipo incorpore el proxy final por intuición.
\end{itemize}

\begin{decision}
SIAP se conserva como fuente pública documentada y señal de stress. No recibe peso positivo
obligatorio en 04F.
\end{decision}
